\chapter{Standards, industries and financing instruments}
\label{ch:modules}

\section{A common core with distinct accounting books}

The common core governs evidence, time, money, events, journals and provenance.
Reporting standards govern recognition, measurement and presentation within
that core. The design preserves three levels of meaning: economic terms such
as a promise to pay cash; an accounting treatment under a particular standard;
and a company's reported label or extension. These levels are connected by
explicit rules and mappings, rather than collapsed into one universal chart
of accounts.

A rule package includes a standards register with jurisdiction, issuer class,
effective dates, transition provisions, permitted elections, paragraph
references, implementation version and tests. The package selector uses the
entity, reporting basis, period and adopted elections. It must distinguish a
standard's mandatory effective date from permitted early adoption. For example,
the IFRS Foundation's IFRS 18 overview states an effective date of 1 January
2027 with earlier application permitted \citep{ifrs18}. A 2026 historical
report cannot be assigned that presentation regime solely because the latest
installed package includes it.

IFRS, US GAAP and ASBE packages share economic contracts and events where
appropriate, but produce separate accounting books. A reconciliation between
books consists of explained adjustments with origins in differing treatments,
not a renaming of accounts. The first implementation can preserve all three
source bases even while only a selected subset has executable rules. Such a
source is represented as observed information with unsupported simulation
treatments, rather than being silently converted to the supported basis.

The official IFRS overviews for IFRS 9, IAS 32 and IFRS 17 establish the relevant
scope of financial-instrument measurement, presentation and insurance
contracts \citep{ifrs9,ias32,ifrs17}. They are orientation sources, not substitutes
for reviewed paragraph-level implementations. Public retrieval of the selected
Chinese Ministry of Finance page and SEC API documentation failed during this
campaign; the source register records these access gaps. ASBE and US GAAP
semantic modules therefore have explicit primary-text acquisition and
accounting-review tasks before they can be marked executable and compliant.

\section{Non-financial companies}

The initial industry package covers sales, collections, purchases, payments,
inventory movements, payroll, capital expenditure, depreciation and basic
financing. Revenue recognition is a separate rule from invoicing and cash
collection. A service contract may generate a performance event before an
invoice, creating a different asset from an unconditional receivable. Customer
advances create obligations before revenue. Returns, rebates and variable
consideration need separate assumptions and subsequent adjustments.

Inventory state includes quantity, cost, location, ownership and lot or cohort.
A costing method is selected within the reporting package's permissible
domain. An aggregate stock can support a simple model only if its approximation
is declared. A commodity-price shock can expose differences between old and
new inventory that an average value suppresses. Fixed-asset state likewise
requires acquisition dates, components, useful lives, residual values,
depreciation method, tax basis and impairment history when those affect the
requested forecast.

Industry extensions add economic mechanisms without bypassing the accounting
core. Construction requires performance and completion states; extractive
industries require reserves, depletion and restoration obligations; utilities
may require rate-regulation treatment; software businesses may require
contract liabilities and development-cost decisions. The module manifest
states what is represented, what is executable, and what remains aggregated.

\section{Debt and hybrid financing}

The contract layer for a debt instrument stores issuer and holder perspective,
currency, issue and settlement dates, issue proceeds and legal face amount.
Its payment terms include coupon/index/spread, reset and fixing dates,
day-count convention, payment calendar, amortisation schedule and maturity.
Its rights and protections include seniority, collateral, covenants, options,
conversion terms, guarantees and linked hedges. Missing terms are
explicit unknowns. A single carrying amount cannot stand in for this record.

The accounting layer stores the classification decision and its evidence,
initial measurement, carrying amount, accrued interest, effective-rate method,
unamortised fees or discounts, fair-value disclosures, impairment or credit
adjustments where applicable, and presentation mappings. Contractual interest,
effective-interest expense, realised cash payments and fair-value changes are
separate quantities. The zero-coupon example in Chapter~\ref{ch:execution}
provides a minimal conformance case for this separation.

The lifecycle must support issuance, recognition, accrual, reset, coupon,
principal repayment, modification, derecognition, default, recovery and
conversion. A modification event first determines which accounting operation
the selected standard requires; it does not always mean issuing a new loan.
An embedded option can affect valuation and classification without being
exercised. A derivative linked as a hedge requires both the economic contract
and any applicable accounting designation; economic risk reduction alone does
not establish hedge-accounting eligibility.

Hybrid instruments make interpretation particularly consequential. IAS 32's
overview discusses liability/equity classification and the separation of a
compound instrument \citep{ias32}. The executable classification must consider
the full terms and the selected standard version. US GAAP or ASBE treatment
must be implemented under its own rules rather than inferred from an IFRS
classification. Where the terms are not disclosed sufficiently, alternative
classifications can coexist only if each is actually compatible with the
known terms and applicable rules.

\section{Equity and ownership}

Equity has both monetary accounts and security quantities. The security
register records share class, issued and outstanding quantities, treasury
holdings, par or stated value where applicable, voting rights, dividend rights,
redemption features, options and conversion relationships. Monetary equity
accounts include contributed capital, reserves, retained earnings and other
components required by the selected basis. These records evolve together
through issuance, repurchase, cancellation, split, dividend and conversion
events.

A stock split changes quantities and per-share measures without necessarily
changing aggregate monetary equity. A cash dividend affects distributions and
cash rather than current operating expense under the corresponding ordinary
equity treatment. Repurchases and employee share awards require their own
rules. Earnings per share uses a dated quantity schedule and rules for potential
ordinary shares; simply dividing profit by the closing share count is
insufficient when shares change during the period.

Funding simulations must connect the quantity register to market assumptions.
An equity issue at a scenario-dependent price changes cash raised, shares and
dilution. The model must state whether it targets a fixed cash amount, a fixed
share count or a capital ratio. These are different management decisions even
when the historical statement would report the same aggregate proceeds.

\section{Banks: accounting capital is not regulatory capital}

A bank package needs instrument and cohort states for loans, securities,
deposits, wholesale funding, repos, derivatives, collateral and commitments.
For interest-rate analysis, currency and repricing buckets matter alongside
remaining maturity. For credit analysis, exposure, default state, loss severity,
recoveries, guarantees and their dependencies matter. For liquidity analysis,
deposit behaviour, collateral eligibility, encumbrance and payment timing matter.
These states can be represented at an explicitly chosen aggregation level;
the public record rarely identifies every contract.

Loan origination can debit loans and credit a customer deposit in the bank's
book. Subsequent transfer to another bank produces settlement flows. Treating
all loan growth as an immediate cash outflow would therefore encode a specific
funding assumption, not a universal banking identity. A system-wide model must
also link counterparties and settlement assets; single-company simulation
records external counterparties without claiming to close the whole economy.

An allowance calculation is a dated measurement based on forecast information.
For a single horizon, let EAD be a fixed exposure at default, PD the probability
of default, and LGD the mean loss fraction conditional on default. Ignoring
discounting, the illustrative expected loss in currency units is
\[
 L=EAD\times PD\times LGD.
\]
With exposure 100, default probability 0.02 and loss given default 0.4, it
produces 0.8. Changing those inputs to 0.04 and 0.5 produces 2, so the allowance
increase is 1.2 if there are no other movements. This calculation is a
demonstration model. A standard-specific expected-credit-loss implementation
must determine horizon, staging or other classification, discounting,
forward-looking information, write-offs and purchased-credit-impaired or
other special cases. The simple formula is not an IFRS 9 or US GAAP compliance
implementation.

Regulatory capital is computed in a separate prudential layer with its own
perimeter, eligibility, deductions, risk weights and transition provisions.
Accounting equity feeds a reconciliation to capital; it is not automatically
Common Equity Tier 1. Liquidity and capital constraints influence the
management policy, which in turn generates asset sales, funding or capital
events. A breached ratio can be an economically meaningful scenario outcome.
No hidden capital injection is permitted to repair it.

\section{Insurance: cohorts and future service}

An insurer requires policy cohorts or groups with issue dates, coverage,
premiums, benefits, expenses, lapses, claims development, guarantees and
reinsurance relationships. Asset portfolios need their own contracts and
valuations. The liability state must distinguish expected future cash flows,
discounting, risk adjustments or margins, unearned profit measures where
applicable, incurred claims and movements between opening and closing states.
The IFRS 17 overview's distinction between measuring future cash flows and
recognising profit as service is provided motivates this separation
\citep{ifrs17}.

The module API accepts policy/cohort state, economic assumptions, non-economic
assumptions and information date. It returns a measurement vector, movement
attribution and proposed recognition events. Experience differences, changes
in assumptions, finance effects and service effects remain separately
identifiable. Reinsurance held has its own record and measurement; it is not
an automatic percentage offset to every insurance liability.

Scenarios can alter discount curves, mortality, morbidity, lapse rates,
inflation, catastrophe losses and policyholder behaviour. Their joint
dependence and cohort exposure must be declared. A portfolio total can match
history while hiding very different guarantee and duration profiles. This
makes alternative economic histories especially valuable, but also limits
what public-data reconstruction can identify. A full insurance package needs
actuarial and accounting review beyond the core release.

\clearpage
\section{Capability contracts for sector packages}

\begin{longtable}{L{26mm}L{55mm}L{55mm}}
\toprule Package & Required state and operations & Minimum acceptance evidence\\\midrule\endhead
Core corporate & Working-capital bridges, lots/assets, ordinary debt and equity,
period closing and cash-flow mappings & Exact synthetic cases, reviewed public
filing reconstruction, all statement reconciliations.\\
Financial instruments & Terms, calendars, resets, options, measurement
categories, modifications and hedges & Contract schedule comparisons, valuation
checks, paragraph-based posting cases and transition cases.\\
Banking & Credit cohorts, deposits, collateral, liquidity, regulatory capital,
management actions & Allowance roll-forwards, repricing and runoff stresses,
capital reconciliation and default/funding-failure cases.\\
Insurance & Policy groups, future cash flows, service and finance components,
reinsurance and asset interaction & Cohort movement reconciliation, assumption
change cases, service recognition and reviewed actuarial calculations.\\
Multi-basis reporting & Versioned IFRS/US GAAP/ASBE rules, entity elections and
transition treatment & Matched economic-event cases with explained basis
differences and independent accounting review.\\\bottomrule
\end{longtable}

Every package declares capabilities at operation level. ``Supports bonds'' is
too vague; ``supports fixed-rate, single-currency, non-callable amortised-cost
issuance and redemption under this reviewed rule set'' is implementable and
testable. Unimplemented operations retain their source and contract
representation while producing an explicit capability error when execution
requires them.
